% Generated by tools/render_documents.py from project/document_projection_map.json.
% Do not edit directly; update the structured projection map and rerun the renderer.

\newcommand{\DoctoralAgendaScopePT}{Esta agenda reúne extensões prospectivas derivadas das limitações do estudo de mestrado. Ela não constitui tese, projeto doutoral formal, aprovação institucional ou compromisso de execução.}
\newcommand{\DoctoralFutureQuestionsPT}{Questões futuras incluem generalização para múltiplas falhas e processos, avaliação de arquiteturas locais distintas, diagnóstico nominal de falhas não vistas, integração controlada entre evidência estatística e explicação por LLM, robustez a mudanças de domínio, validação com dados industriais autorizados e estudo de governança para ambientes regulados.}
\newcommand{\LiteratureStatusEN}{The literature review remains under construction. The Xing, Khan, and Kai triad is a preliminary basis, but bibliographic identity, legitimate full text, precise locators, relevance to IDV(13), and claim-support capacity still require verification. The gap matrix must include zero-shot detection, normal-only detection, and direct LLM time-series antecedents before any novelty claim is made.}
\newcommand{\LiteratureStatusPT}{A revisão bibliográfica está em construção. A tríade Xing, Khan e Kai constitui uma base preliminar, mas identidade bibliográfica, texto integral, locadores, aderência à IDV(13) e capacidade de sustentar alegações ainda precisam ser verificadas. A matriz de lacuna deve incluir antecedentes de detecção zero-shot, normal-only e uso direto de LLM em séries temporais antes de qualquer alegação de novidade.}
\newcommand{\ResearchContributionEN}{The proposed contribution is a reproducible protocol for evaluating fault-zero-shot anomaly detection by general-purpose local LLMs on multivariate industrial temporal data, including normal-operation reference, causal windows, alarm and abstention rules, evidence-coherence evaluation, and verifiable local execution.}
\newcommand{\ResearchContributionPT}{A contribuição proposta é um protocolo reproduzível para avaliar detecção zero-shot de anormalidades por LLMs de propósito geral executadas localmente sobre dados temporais industriais multivariados, incluindo referência de normalidade, janelas causais, regras de alarme e abstenção, avaliação de coerência das evidências e execução local verificável.}
\newcommand{\ResearchHypothesisEN}{H1: a general-purpose LLM contextualized only with process knowledge and normal operation can detect an abnormal condition caused by a previously unseen fault without fault-specific training. H2: detection capability increases as multivariate evidence accumulates across causal temporal windows, producing an observable transition from NORMAL or EVIDENCE_INSUFFICIENT to ANOMALY after true fault onset. H3: when the LLM indicates an anomaly, the variables and trends cited in its justification show verifiable correspondence with changes actually observable in the system.}
\newcommand{\ResearchHypothesisPT}{H1: uma LLM de propósito geral contextualizada apenas com o processo e a operação normal é capaz de detectar uma condição anormal provocada por uma falha não vista, sem treinamento específico para a falha. H2: essa capacidade de detecção aumenta à medida que evidências multivariadas se acumulam em janelas temporais causais, com transição observável de NORMAL ou EVIDÊNCIA_INSUFICIENTE para ANOMALIA após o início real da falha. H3: quando a LLM indica anomalia, as variáveis e tendências citadas em sua justificativa apresentam correspondência verificável com alterações efetivamente observadas no sistema.}
\newcommand{\ResearchLimitationsEN}{The TEP is public and simulated; therefore, the experiment may demonstrate a non-externalizing architecture but cannot prove protection of real industrial secrets or complete compliance. Local execution is a measurable experimental condition rather than a sufficient novelty claim. The specific gap concerning normal-reference fault-zero-shot detection in multivariate industrial systems remains provisional pending deeper literature review.}
\newcommand{\ResearchLimitationsPT}{O TEP é público e simulado; portanto, o experimento poderá demonstrar uma arquitetura sem externalização, mas não provar proteção de segredos industriais reais ou conformidade completa. A execução local é condição experimental mensurável, não novidade suficiente. A lacuna específica de detecção zero-shot normal-reference em sistemas industriais multivariados permanece provisória até revisão bibliográfica aprofundada.}
\newcommand{\ResearchMethodEN}{The study uses the Tennessee Eastman Process as a public simulated environment, provisionally focusing on IDV(13). Simulator-controlled fault activation provides the experimental ground truth. DPCA with T-squared and SPE/Q serves as an independent statistical reference. The LLM is executed locally without an external API and without fault-specific training, labeled IDV(13) examples, or the IDV(13) description in the primary condition; it may receive process knowledge, normal-operation reference information, and a structured temporal representation through causal windows.}
\newcommand{\ResearchMethodPT}{O estudo utilizará o Tennessee Eastman Process como ambiente público e simulado, com foco provisório na IDV(13). O instante de ativação controlado pelo simulador fornecerá o ground truth experimental. A DPCA com T2 e SPE/Q será uma referência estatística independente. A LLM será executada localmente, sem API externa e sem treinamento específico, exemplos rotulados ou descrição da IDV(13) na condição primária; poderá receber conhecimento do processo, referência de operação normal e representação temporal estruturada em janelas causais.}
\newcommand{\ResearchObjectiveEN}{To evaluate, through progressive causal windows of the Tennessee Eastman Process, the ability of a local fault-blind LLM to detect the abnormal condition induced by IDV(13), using the true simulator-controlled activation time as experimental ground truth and DPCA as an independent statistical reference.}
\newcommand{\ResearchObjectivePT}{Avaliar, em janelas temporais progressivas e causais do Tennessee Eastman Process, a capacidade de uma LLM local e fault-blind de detectar o surgimento da anormalidade provocada pela IDV(13), usando o instante real de ativação como ground truth experimental e DPCA como referência estatística independente.}
\newcommand{\ResearchProblemEN}{This study investigates whether a general-purpose LLM, executed locally and contextualized only with process knowledge and normal operation, can detect the emergence of an abnormal condition in a multivariate industrial system under a fault-zero-shot protocol.}
\newcommand{\ResearchProblemPT}{Investigar se uma LLM de propósito geral, executada localmente e contextualizada apenas com conhecimento do processo e de sua operação normal, é capaz de detectar o surgimento de uma condição anormal em um sistema industrial multivariado em regime zero-shot quanto à falha.}
\newcommand{\ResearchQuestionEN}{Can a general-purpose LLM, without fault-specific training and contextualized only with knowledge of the process and its normal operation, detect in a fault-zero-shot setting the emergence of an abnormal condition in a multivariate industrial system?}
\newcommand{\ResearchQuestionPT}{Uma LLM de propósito geral, sem treinamento específico para falhas e contextualizada apenas pelo conhecimento do processo e de sua operação normal, é capaz de detectar em regime zero-shot o surgimento de uma condição anormal em um sistema industrial multivariado?}
\newcommand{\ResearchTitleEN}{Zero-Shot Fault Detection by a General-Purpose LLM in a Multivariate Industrial System: A Tennessee Eastman Process Study with IDV(13)}
\newcommand{\ResearchTitlePT}{Detecção zero-shot de falhas por LLM em sistema industrial multivariado: estudo no Tennessee Eastman Process com IDV(13)}
\newcommand{\ResultsStatusEN}{No experimental results have yet been accepted. This version records the design and hypotheses H1-H3; results may be incorporated only after reproducible execution, artifact preservation, metric calculation, and evaluation of a frozen snapshot.}
\newcommand{\ResultsStatusPT}{Ainda não existem resultados experimentais aceitos. Esta versão registra o desenho e as hipóteses H1-H3; resultados somente poderão ser incorporados após execução reproduzível, preservação dos artefatos, cálculo das métricas e avaliação do snapshot congelado.}
